%% Career-Ops LaTeX CV Template
%% Style: ATS-optimized single-column, minimal color accent, single-page friendly.
%% Build: pdflatex {file}.tex (or tectonic {file}.tex)

\documentclass[letterpaper,10pt]{article}

\usepackage[margin=0.55in]{geometry}
\usepackage[T1]{fontenc}
\usepackage[utf8]{inputenc}
\usepackage{lmodern}
\usepackage{microtype}
\usepackage{enumitem}
\usepackage{titlesec}
\usepackage{hyperref}
\usepackage{xcolor}
\usepackage{parskip}

\definecolor{accent}{HTML}{1F3A5F}

\hypersetup{
  colorlinks=true,
  urlcolor=accent,
  linkcolor=accent,
  pdftitle={ Ahnaf An Nafee - Resume },
  pdfauthor={ Ahnaf An Nafee },
  pdfsubject={Resume},
  pdfkeywords={ Release Engineering, CI/CD, Kubernetes, OpenShift, AWS, Go, Java, Python, Reproducible Builds, Packaging, Container Orchestration, Helm, Gradle, Maven, Bazel, Observability, AI Infrastructure, On-Prem Deployment, Self-Hosted, Distributed Systems, Backend Systems, Infrastructure Automation, RBAC }
}

% Ensure generated PDF is machine readable / ATS parseable
\input{glyphtounicode}
\pdfgentounicode=1

\pagestyle{empty}
\setlength{\parindent}{0pt}
\raggedbottom
\raggedright

% --- Section formatting ------------------------------------------------------
% Note: section titles render in mixed case (no \MakeUppercase) for maximum
% ATS compatibility. Older parsers normalize headers like "Work Experience"
% and "Education" more reliably than uppercase variants.
\titleformat{\section}
  {\color{accent}\large\bfseries}
  {}{0em}{}[{\color{accent}\vspace{-6pt}\hrulefill\vspace{-4pt}}]
\titlespacing*{\section}{0pt}{8pt}{4pt}

\titleformat{\subsection}
  {\normalsize\bfseries}
  {}{0em}{}
\titlespacing*{\subsection}{0pt}{4pt}{2pt}

% --- Custom commands ---------------------------------------------------------
\newcommand{\role}[4]{%
  \textbf{#1} \hfill \textit{#3}\\[-1pt]
  \textit{#2} \hfill \textit{#4}\par\vspace{2pt}
}

\newcommand{\education}[4]{%
  \textbf{#1} \hfill \textit{#3}\\[-1pt]
  \textit{#2} \hfill \textit{#4}\par\vspace{2pt}
}

\newlist{tightitemize}{itemize}{1}
\setlist[tightitemize]{leftmargin=1.3em, topsep=2pt, itemsep=1pt, parsep=0pt, label=\textbullet}

% =============================================================================
\begin{document}

% --- Header ------------------------------------------------------------------
\begin{center}
  {\LARGE\bfseries\color{accent} Ahnaf An Nafee }\\[3pt]
  {\normalsize Software Engineer \textbar{} AI Researcher \textbar{} DevOps \& Cloud Infrastructure }\\[4pt]
  \small
  \href{mailto:ahnafnafee@gmail.com}{ahnafnafee@gmail.com} \textbar{}
  540-252-8738 \textbar{}
  Fairfax, VA\\
  \href{https://linkedin.com/in/ahnafnafee}{linkedin.com/in/ahnafnafee} \textbar{}
  \href{https://github.com/ahnafnafee}{github.com/ahnafnafee} \textbar{}
  \href{https://www.ahnafnafee.dev/}{ahnafnafee.dev} \textbar{} \href{https://scholar.google.com/citations?user=u15DO0cAAAAJ}{Google Scholar}
\end{center}

% --- Summary -----------------------------------------------------------------
\section*{Summary}
DevOps and Release Engineer with 4+ years building reproducible CI/CD pipelines, Kubernetes/OpenShift platforms, and multi-language container packaging across Go, Java, and Python services. Proven release-engineering record at Paychex: cross-team Kubernetes standardization, automated TLS lifecycle eliminating 100\% of certificate-driven downtime, and a Gradle plugin for image certification that cut build congestion and 10\% annual opex. Former CTO at Dynasty 11 — led AWS migration with 80\% reduction in load and cost, and shipped GitHub Actions CI/CD reducing manual release effort by 85\%. Active CS PhD researcher at George Mason DCXR Lab whose AI/3D-graphics work directly overlaps Arize's AI observability mission. Seeking a senior DevOps / Release Engineering role packaging and shipping AI infrastructure into self-hosted, on-prem customer environments.

% --- Skills ------------------------------------------------------------------
\section*{Technical Skills}
\textbf{Release \& Build:} CI/CD, GitHub Actions, Jenkins, Gradle, Maven, Helm, Reproducible Builds, Container Image Certification, Packaging, Release Pipelines\\
\textbf{Cloud \& Orchestration:} AWS, Kubernetes, OpenShift, Docker, Terraform, RBAC, On-Prem / Self-Hosted Deployment Patterns, Observability, High Availability\\
\textbf{Languages:} Python, Go, Java, C++, TypeScript, JavaScript, SQL, GLSL, Bash\\
\textbf{Backend \& Systems:} Distributed Systems, Microservices, REST, gRPC, WebSockets, OAuth, System Design, Infrastructure-as-Code\\
\textbf{AI / ML:} PyTorch, TensorFlow, Generative AI, Deep Learning, Computer Vision, Diffusion Models, LLMs, MLOps\\
\textbf{Graphics \& XR:} 3D Computer Graphics, Real-Time Rendering, Shader Programming, Unity, Unreal Engine, WebGL, Procedural Generation\\
\textbf{Methods:} Agile, Scrum, Test-Driven Development, Code Review, Technical Leadership, Cross-Functional Collaboration

% --- Experience --------------------------------------------------------------
\section*{Work Experience}
\role{DevOps Engineer}{Paychex}{Feb 2023 -- Aug 2025}{Rochester, NY}
\begin{tightitemize}
  \item Built a Gradle plugin for container image certification that standardized reproducible build pipelines across enterprise services, reduced CI/CD pipeline congestion, and lowered operational expenses by \textbf{10\% annually}
  \item Architected and led cross-team standardization of Kubernetes/OpenShift deployment dictionaries across 10+ service teams, eliminating configuration conflicts and accelerating release velocity
  \item Automated end-to-end TLS certificate lifecycle management, eliminating manual renewals and reducing certificate-expiration-related downtime by \textbf{100\%} — improving platform reliability and upgradeability
  \item Designed and rolled out Role-Based Access Control (RBAC) policies in OpenShift in partnership with security and infrastructure teams, hardening the platform and enabling enterprise-grade, air-gapped-friendly deployment patterns
  \item Championed infrastructure-as-code and observability best practices; partnered with SREs to improve deployment reliability, MTTR, and release cadence
\end{tightitemize}

\role{Chief Technology Officer}{Dynasty 11 Studios}{Jun 2022 -- Feb 2023}{Wayne, PA}
\begin{tightitemize}
  \item Introduced CI/CD with GitHub Actions and Maven, automating builds, tests, and deployments across Java/Go services; reduced manual release effort by \textbf{85\%} and eliminated deploy-day toil
  \item Spearheaded migration of monolithic backend to microservices on AWS, cutting application load and cloud spend by \textbf{80\%} through right-sizing, autoscaling, and serverless adoption
  \item Shipped reproducible builds and packaging conventions across Java/Go services, enabling consistent multi-environment delivery
  \item Owned end-to-end technical strategy, architecture, and execution; led engineering across backend, frontend, and cloud teams shipping a consumer-facing matchmaking platform
  \item Recruited, mentored, and led multiple engineering pods through full microservice design, implementation, and production rollout; ran weekly stakeholder syncs aligning product, design, and engineering
\end{tightitemize}

\role{Software Engineer}{Dynasty 11 Studios}{Sep 2021 -- Jun 2022}{Wayne, PA}
\begin{tightitemize}
  \item Engineered a real-time chat subsystem using STOMP over WebSockets in Java, supporting thousands of concurrent users with low-latency delivery
  \item Integrated 20+ REST endpoints across third-party and OAuth-based identity providers, enabling secure, scalable authentication and data exchange
  \item Optimized client-side state management with Redux, cutting chat-service performance bottlenecks by \textbf{80\%}
  \item Designed data pipelines and APIs powering the matchmaking recommendation engine serving thousands of active users
\end{tightitemize}

\role{Graduate Teaching Assistant}{George Mason University --- DCXR Lab}{Aug 2025 -- Present}{Fairfax, VA}
\begin{tightitemize}
  \item Mentor 60+ graduate and undergraduate students in algorithms, distributed systems, and applied machine learning, driving measurable improvements in lab completion and conceptual mastery
  \item Design coursework, lab exercises, and evaluation rubrics aligned with modern software engineering and infrastructure best practices
\end{tightitemize}

\role{Technical Programmer}{PHL Collective}{Mar 2021 -- Sep 2021}{Philadelphia, PA}
\begin{tightitemize}
  \item Shipped gameplay systems and parameterized shaders for \textit{DC's Justice League: Cosmic Chaos} (multi-platform console/PC release)
  \item Built modular game-manager scripts that streamlined designer workflows and accelerated feature iteration
  \item Partnered with technical artists to create customizable, performance-tuned shaders optimized for runtime rendering
  \item Executed integration and stress testing via Mantis, delivering actionable, developer-ready bug reports
\end{tightitemize}

% --- Education ---------------------------------------------------------------
\section*{Education}
\education{Ph.D. in Computer Science}{George Mason University --- DCXR Lab}{}{Fairfax, VA}
\begin{tightitemize}
  \item Research focus: generative AI for 3D content creation, machine-learning-driven graphics pipelines (UV mapping, non-photorealistic rendering), and human--computer interaction in immersive XR environments — directly relevant to AI observability and infrastructure for ML systems
\end{tightitemize}

\education{B.S. in Computer Science}{Drexel University}{}{Philadelphia, PA}
\begin{tightitemize}
  \item Honors: Dean's List; Founder's Scholarship
  \item Relevant coursework: Algorithms, Operating Systems, Distributed Computing, Computer Graphics, Machine Learning, Software Engineering
\end{tightitemize}

% --- Projects ----------------------------------------------------------------
% Section header is "Selected Projects" (not "Research & Selected Projects")
% so ATS classify the items below as Projects. Research focus stays in the
% bullet text itself.
\section*{Selected Projects}
\begin{tightitemize}
\item \textbf{AI-Driven 3D Content Generation} --- Researching generative models (diffusion, neural fields) for automated, controllable 3D asset creation in production-grade graphics pipelines
\item \textbf{ML for Graphics Pipelines} --- Exploring learning-based approaches to UV mapping, stylization, and non-photorealistic rendering (NPR) for real-time applications
\item \textbf{Portfolio \& Publications} --- Additional projects, technical artwork, and game development work available at \href{https://www.ahnafnafee.dev/portfolio}{ahnafnafee.dev/portfolio}
\end{tightitemize}

\end{document}
